\section{The later supplements, 5--9 September}\label{sec:later}

The two later versions of the record add work that uses the construction's data as stated input. Both say that they do not by themselves establish the construction or a counterexample to \cite{CDP20} (guide \fn{14\_}, §§1 and~4; key-advances reader \fn{26\_}, closing paragraph). This section maps them, reproduces their finite statements, and proves one arithmetic statement in full.

\subsection{The topology supplement of 5 September (version 1.2)}

The supplement consists of two papers, \emph{The higher torus wrap} (\fn{15\_}, 50 pages) and \emph{Real $S^{12}$ and $S^{24}$ routes} (\fn{16\_}, 64 pages), and a three-page reader's guide (\fn{14\_}). All three concern maps out of the cusp region of $X$:
\begin{enumerate}[label=\arabic*.]
\item \emph{The angle-forgetting map.} A map $g:X\to S^3$ built from the cusp quotient is based null-homotopic. The proof exhibits an explicit framed null-bordism of the inverse image, a normally framed $T^3$, inside one cusp chart times an interval. It does not use any identification of $X$ with a sphere.
\item \emph{The retained-angle map.} For $H:X\to S^4$ and the reconstruction's fixed diffeomorphism $d:S^6\to X$, the record proves $[H\circ d]=\eta_4\circ\eta_5\neq0$ in $\pi_6(S^4)\cong\ZZ/2$. The inverse image is a framed $T^2$ with the Lie framing. Its spin quadratic form $q(x,y)=x+xy+y$ over $\Ff_2$ has Arf invariant $1$, since $q$ takes the value $1$ on three of the four vectors. This uses the global sphere identification as input.
\item \emph{A Hopf morphism.} Smashing with the Hopf map, $x\mapsto\eta_2\wedge x$, gives an injective homomorphism $\pi_6(S^4)\to\pi_9(S^6)\cong\ZZ/24$ sending $[H\circ d]$ to $12\nu_6$. This is the classical relation $\eta^3=12\nu$ \cite{Toda62} evaluated on the class of item~2, as the guide itself says.
\item \emph{A peripheral morphism.} With the peripheral character $\chi:SL_2(\ZZ)\to\ZZ/12$ and the values $(\chi(C_0),\chi(C_1),\chi(C_2'))=(-1,4,-3)$, the map $\Psi_H=12\chi\bmod24$ takes the values $(12\nu_6,0,12\nu_6)$, and its kernel is $\chi^{-1}(2\ZZ/12)$.
\item \emph{Higher tori.} Over the cusp there is a local map of complex three-tori onto the two-tori, with elliptic kernel, which extends over the toric fillings. The papers do not construct a global higher-torus family. The direct rank-three cusp would have Euler characteristic $6$, and with two Euler-zero finite fibres the total would be $6$, not $\chi(S^8)=2$.
\item \emph{Octonionic cones.} The equal-diagonal positive cone of the octonionic Jordan algebra, compactified at one point, is the pair $(B^{25},S^{24})$. The papers relate it to the flag manifold $F_4/\mathrm{Spin}(8)$; neither completes an analytic torus gluing.
\end{enumerate}
\status{Items 1--3, 5 (the local map) and 6: \textsf{located}. \textsf{Reproduced}: the Arf invariant in item~2; the values and kernel in item~4; the Euler-characteristic arithmetic in item~5. The groups $\pi_6(S^4)\cong\ZZ/2$ (generated by $\eta_4\eta_5$) and $\pi_9(S^6)\cong\ZZ/24$, and $\eta^3=12\nu$, are classical \cite{Toda62}.}

\subsection{The key-advances reader of 6 September}

The five-page reader \fn{26\_} gives formulas and proof locations for five results. The proofs are in the frozen project archive (\fn{28\_}; its 782-page workbench and 138-page higher-rung paper), which was not downloaded here.
\begin{enumerate}[label=\textsf{S6-\arabic*}, leftmargin=3.4em]
\item \emph{The anticanonical ring recovers the fibration.} With $f^*(p_1)=3S_1$ and $f^*(p_2)=4S_2$, the reader states $\omega_X\cong\mathcal O_X(-2S_2)$, $\omega_X^{-2}\cong f^*\mathcal O_B(p_2)$ and $f_*\mathcal O_X(kS_2)=\mathcal O_B(\lfloor k/4\rfloor p_2)$. Hence $\bigoplus_mH^0(X,\omega_X^{-m})\cong\C[U,V]$ with $\deg U=1$ and $\deg V=2$, and the degree-two part recovers $f$.
\item \emph{A nonzero deformation.} Replacing $\beta$ by $\beta+c$, with $\operatorname{Im}c$ below an explicit bound, gives a proper holomorphic submersion over a disc of values of $c$. Its Kodaira--Spencer class at each point is nonzero.
\item \emph{The finite fillings.} There are explicit product covers of degrees nine and eight of the central surfaces of the two multiple fibres, and the kernel of the attachment map is an infinite cyclic central subgroup.
\item \emph{Niemeier--Leech arithmetic.} Take the Niemeier lattice $N$ with root system $A_5^4D_4$ and glue index $72$. A neighbour construction with the Weyl vector gives three neighbours, which are cycled by an order-three action. Their common intersection $J$ has determinant $81$ and index $9$ in each, and the fixed extension has root system $A_2^{12}$ with $72$ roots. The value groups of an octonionic cubic on these lattices are computed, with inclusions of index $4$ and $3$.
\item \emph{A common lattice in the triality space.} In Wilson's octonionic Leech construction \cite{Wilson09}, the intersection of the Leech copy with the transported Niemeier lattices is $I=\{\sqrt2(a,a,a):a\in D_4\}$, with Gram matrix $6G_{D_4}$, determinant $5184$ and minimum $12$. On it the cubic is $\tau(\sqrt2(a,a,a))=2\sqrt2\,P(a)$ with $P(a)=a_0\bigl(a_0^2-3(a_1^2+a_2^2+a_3^2)\bigr)$. It generates the group $4\sqrt2\,\ZZ$ but does not take the value $12\sqrt2$.
\end{enumerate}

\begin{proposition}[the arithmetic of S6-5]\label{prop:s65}
Let $D_4=\{a\in\ZZ^4:a_0+a_1+a_2+a_3\equiv0\pmod2\}$, let $a$ run over $D_4$, viewed as quaternions $a_0+a_1i+a_2j+a_3k$, and put $\tau(q)=\operatorname{Re}\bigl((qq)q\bigr)$.
\begin{enumerate}[label=\textup{(\alph*)}]
\item $\tau(\sqrt2\,a)=2\sqrt2\,P(a)$.
\item $P(a)$ is even for every $a\in D_4$, and $P(1,1,0,0)=-2$. So the values $\tau(\sqrt2\,a)$ generate $4\sqrt2\,\ZZ$.
\item For every $a\in\ZZ^4$, if $3\mid P(a)$ then $9\mid P(a)$. So $P$ never takes a value $3k$ with $3\nmid k$. In particular, since $\tau(\sqrt2\,a)=12\sqrt2\,k$ means $P(a)=6k$, $\tau(\sqrt2\,a)\neq12\sqrt2\,k$ whenever $3\nmid k$; the key-advances reader states the case $k=1$.
\end{enumerate}
\end{proposition}

\begin{proof}
(a) Write $a=a_0+v$ with $v$ purely imaginary, so $v^2=-|v|^2$. Then $a^2=a_0^2-|v|^2+2a_0v$, and
\[
\operatorname{Re}(a^3)=a_0(a_0^2-|v|^2)+2a_0\operatorname{Re}(v^2)=a_0(a_0^2-3|v|^2).
\]
Multiply by $(\sqrt2)^3=2\sqrt2$. The quaternions form a subalgebra of the octonions, so the octonionic product agrees.

(b) If $a_0$ is even, so is $P(a)$. If $a_0$ is odd, then $a_1+a_2+a_3$ is odd, so $|v|^2\equiv a_1+a_2+a_3\equiv1$ and $a_0^2-3|v|^2\equiv1-3\equiv0\pmod2$.

(c) $P(a)\equiv a_0^3\pmod3$. So $3\mid P(a)$ forces $a_0=3b$, and then $P(a)=9b(3b^2-|v|^2)$.
\end{proof}

\status{S6-1 to S6-4: \textsf{located}; the proofs are in the archive and were not read. \textsf{Reproduced}:
\begin{itemize}
\item S6-1: the dimension count $h^0(\omega_X^{-m})=h^0(\Pp^1,\mathcal O(\lfloor m/2\rfloor))=\lfloor m/2\rfloor+1$, which matches the graded pieces of $\C[U,V]$. Also $\omega_X\cong\mathcal O_X(-2S_2)$ agrees with the manuscript's $K_X\cong f^*\mathcal O_B(-1)\otimes\mathcal O_X(2S_2)$, because $f^*\mathcal O_B(1)=\mathcal O_X(4S_2)$.
\item S6-2: $\det_{\R}\Pi=\operatorname{Im}\tau\cdot D$, so the determinant ratio $(D+\operatorname{Im}\delta)/D$ holds.
\item S6-4: the index $72=\sqrt{6^4\cdot4}$; the $72$ roots of $A_2^{12}$; the determinant $81=9^2$.
\item S6-5: the Gram determinant, the minimum, and the formula for $\tau$.
\end{itemize}
Proposition~\ref{prop:s65}: \textsf{proved here}. Parts (b) and (c) are also in the reader of four workbenches \cite[Prop.~4.7]{WorkbenchReader}, which took the formula for $\tau$ from the source; part (a) checks that formula. The mod-$3$ step of (c) is the key-advances reader's own argument.}
